The two branches share only the deterministic external inputs in the statement; their law classes are distinct.

Along the stated sequence of non-Gaussian experiments, \(\text{thm:adaptive-rootn-minimax}\) supplies constants \(0<c<C<\infty\), independent of \(n\), such that
\[
 \frac{c}{n}
 \le \mathfrak R_n^{\mathrm{NG}}
 \le \frac{C}{n}.
\]
Therefore \(\mathfrak R_n^{\mathrm{NG}}\asymp n^{-1}\). The same theorem identifies \(\widehat\theta_{\mathrm{spec},n}\) as the common upper-bound witness and states that it uses \(\bar g_n\), uses no supplied outcome code, and remains valid for fixed nonzero \(\varepsilon_{1,n}\le\varepsilon_0\). The definition of \(\mathcal P_{\mathrm{NG},n}\) likewise contains no restriction involving \(\widehat q_n\), \(\bar q_n\), or \(\varepsilon_{2,n}\). This proves the first branch.

For the second branch, \(\text{prop:bounded-outcome-gaussian-degeneracy}\) proves that every member of \(\mathcal P_{\mathrm G,n}^{\mathrm{JMS}}\) has \(\theta_0(P)=0\). If the class is nonempty, the same proposition applies the zero statistic to the two definitions in \(\text{def:minimax-risks}\) and obtains
\[
 \mathfrak R_n^{\mathrm G}
 =\mathfrak M_{n,1-\gamma}^{\mathrm G}
 =0.
\]
This conclusion is incompatible with any strictly positive lower bound on the displayed intersection. It follows from the simultaneous conditions defining that intersection, not from a transfer of Jin--Mackey--Syrgkanis Theorem 3.2. The proof uses no relation between the two law classes.